\documentclass[10pt,a4paper]{amsart}
\usepackage{amsmath,amssymb,amsthm,graphicx}
% For commutative diagrams we use Xy-pic
\usepackage[all]{xy}

% We use 2cell for 2-commutative diagrams.
\xyoption{2cell}
% Theorem environments.
%
\theoremstyle{plain}
\newtheorem{theorem}[subsection]{Theorem}
\newtheorem{proposition}[subsection]{Proposition}
\newtheorem{lemma}[subsection]{Lemma}

\theoremstyle{definition}
\newtheorem{definition}[subsection]{Definition}
\newtheorem{example}[subsection]{Example}
\newtheorem{exercise}[subsection]{Exercise}
\newtheorem{situation}[subsection]{Situation}

\theoremstyle{remark}
\newtheorem{remark}[subsection]{Remark}
\newtheorem{remarks}[subsection]{Remarks}

\numberwithin{equation}{subsection}

% Macros
%
\def\lim{\mathop{\mathrm{lim}}\nolimits}
\def\colim{\mathop{\mathrm{colim}}\nolimits}
\def\Spec{\mathop{\mathrm{Spec}}}
\def\Hom{\mathop{\mathrm{Hom}}\nolimits}
\def\Ext{\mathop{\mathrm{Ext}}\nolimits}
\def\SheafHom{\mathop{\mathcal{H}\!\mathit{om}}\nolimits}
\def\SheafExt{\mathop{\mathcal{E}\!\mathit{xt}}\nolimits}
\def\Sch{\mathit{Sch}}
\def\Mor{\mathop{\mathrm{Mor}}\nolimits}
\def\Ob{\mathop{\mathrm{Ob}}\nolimits}
\def\Sh{\mathop{\mathit{Sh}}\nolimits}
\def\NL{\mathop{N\!L}\nolimits}
\def\CH{\mathop{\mathrm{CH}}\nolimits}
\def\proetale{{pro\text{-}\acute{e}tale}}
\def\etale{{\acute{e}tale}}
\def\QCoh{\mathit{QCoh}}
\def\Ker{\mathop{\mathrm{Ker}}}
\def\Im{\mathop{\mathrm{Im}}}
\def\Coker{\mathop{\mathrm{Coker}}}
\def\Coim{\mathop{\mathrm{Coim}}}

% Boxtimes
%
\DeclareMathSymbol{\boxtimes}{\mathbin}{AMSa}{"02}

%
% Macros for moduli stacks/spaces
%
\def\QCohstack{\mathcal{QC}\!\mathit{oh}}
\def\Cohstack{\mathcal{C}\!\mathit{oh}}
\def\Spacesstack{\mathcal{S}\!\mathit{paces}}
\def\Quotfunctor{\mathrm{Quot}}
\def\Hilbfunctor{\mathrm{Hilb}}
\def\Curvesstack{\mathcal{C}\!\mathit{urves}}
\def\Polarizedstack{\mathcal{P}\!\mathit{olarized}}
\def\Complexesstack{\mathcal{C}\!\mathit{omplexes}}
% \Pic is the operator that assigns to X its picard group, usage \Pic(X)
% \Picardstack_{X/B} denotes the Picard stack of X over B
% \Picardfunctor_{X/B} denotes the Picard functor of X over B
\def\Pic{\mathop{\mathrm{Pic}}\nolimits}
\def\Picardstack{\mathcal{P}\!\mathit{ic}}
\def\Picardfunctor{\mathrm{Pic}}
\def\Deformationcategory{\mathcal{D}\!\mathit{ef}}
\usepackage{cleveref}
\crefname{originallemma}{original lemma}{original lemmas}
\Crefname{originallemma}{Original lemma}{Original lemmas}
\makeatletter
\newlabel{lemma-silly}{{0F9H}{source}{}{sourcepointer-1}{}}
\newlabel{lemma-silly@cref}{{[originallemma][0][]0F9H}{[1][0][]source}{}{}{}}
\makeatother
\title{Independence of the locally constructed bivariant class from the choice of the complementary closed subscheme}
\author{The Stacks Project Authors}
\date{stacks-project:chow.tex:local-snapshot-20261012; GNU FDL 1.2 or later}
\begin{document}
\maketitle
\noindent\textit{Editorial note.} All delivered mathematical text is the author's own, copied from the pinned Stacks Project chapter file chow.tex; no statement or proof step was rewritten, reordered or removed.\par\medskip
\noindent\textit{Editorial note.} Preamble provenance: the pinned chapter file opens with an input of the shared preamble and carries no preamble text of its own, so the project's single Stacks-wide preamble staged at sources/stacks-chow/preamble.tex is excerpted verbatim as two spans, lines 14-18 (Xy-pic setup) and lines 154-217 (the native theorem declarations and mathematical macros). Nothing was added, removed or redefined; the wrapper class and the amsmath/amssymb/amsthm/graphicx base packages are supplied by the assembler.\par\medskip
\noindent\textit{Editorial note.} The counted statement and proof contain no cross reference, no citation and no included file, so the delivered body is self contained.\par\medskip
\noindent\textit{Editorial note.} Canonical location of the counted proposition: Stacks tag 0FAH (sources/stacks-chow/pinned-original-source.tex line 9752).\par\medskip
\noindent\textit{Allowed standard prerequisites (editorial).}\begin{itemize}
\item The local construction of P'\_p(E\_2) and c'\_p(E\_2) and its two hypotheses. Same chapter, Lemma labelled lemma-silly (source line 9695), whose hypotheses the statement inherits. Fixes the situation X = X\_1 union X\_2, the closed subscheme E\_2 and the condition on the restriction of E\_2 to X\_1 cap X\_2, whose choice is what this lemma shows to be immaterial.
\item Decomposition of a cycle class along a cover by two closed subschemes and bivariant capping. Same chapter, chow.tex, the cycle group Z\_k(X) and the bivariant operation used at proof lines 9764-9769. Supplies that alpha = i\_*beta + i\_\{2,*\}alpha\_2 with beta supported on the intersection and alpha\_2 on X\_2, which is the step that makes the two constructions agree.
\end{itemize}\medskip
\ifdefined\hypertarget\hypertarget{sourcepointer-1}{}\fi
\noindent\textit{Source reference.} Original lemma 0F9H: sources/stacks-chow/pinned-original-source.tex, source lines 9695--9695. Same-chapter pointer into the author environment lemma at source line 9695; invoked at proof line None of the counted proof. It is retained only as a navigation pointer; every load-bearing step is inside the delivered span.\par\medskip
\begin{lemma}
\label{lemma-silly-independent}
In Lemma \ref{lemma-silly} the bivariant class
$P'_p(E_2)$, resp.\ $c'_p(E_2)$ in $A^p(X_2 \to X)$
does not depend on the choice of $X_1$.
\end{lemma}

\begin{proof}
Suppose that $X_1' \subset X$ is another closed subscheme such that
$X = X'_1 \cup X_2$ set theoretically and the restriction
$E_2|_{X'_1 \cap X_2}$ is zero, resp.\ isomorphic to a
finite locally free $\mathcal{O}_{X'_1 \cap X_2}$-module
of rank $< p$ sitting in cohomological degree $0$.
Then $X = (X_1 \cap X'_1) \cup X_2$. Hence we can write
any element $\alpha \in \CH_k(X)$ as $i_*\beta + i_{2, *}\alpha_2$ with
$\alpha_2 \in \CH_k(X'_2)$ and $\beta \in \CH_k(X_1 \cap X'_1)$.
Thus it is clear that
$P'_p(E_2) \cap \alpha = P_p(E_2) \cap \alpha_2 \in \CH_{k - p}(X_2)$,
resp.\  $c'_p(E_2) \cap \alpha = c_p(E_2) \cap \alpha_2 \in \CH_{k - p}(X_2)$,
is independent of whether we use $X_1$ or $X'_1$. Similarly
after any base change.
\end{proof}

\end{document}
